\textbf{P1)} Una partícula realiza un movimiento armónico simple en una dimensión. Su posición está dada por:

\[
x(t) = 0.30\cos\left(2t + \frac{\pi}{6}\right)\ [\mathrm{m}]
\]

donde $t$ se mide en segundos.

\begin{enumerate}[a)]
  \item Determine la amplitud $A$, la frecuencia angular $\omega$, la frecuencia $f$, el período $T$ y la fase inicial $\delta$.
  \item Calcule la posición de la partícula en $t = 1.0\ \mathrm{s}$.
  \item Determine la velocidad $v(t)$ y la aceleración $a(t)$ en cualquier instante.
  \item Calcule la rapidez máxima y la aceleración máxima.
\end{enumerate}

\vspace{0.3cm}